\chapter{Venepuncture}

\section*{Overview}
Venepuncture is the puncturing of a peripheral vein with a needle to obtain a blood sample for laboratory testing. It is one of the most frequently performed procedures in health care, providing blood for hematologic, biochemical, and microbiologic analysis.
\section*{Alternative Names}
Phlebotomy; venipuncture; blood draw; venous blood sampling
\section*{Principle}
A tourniquet distends a superficial peripheral vein, and a beveled needle or vacuum tube system pierces the vein wall to allow blood to flow into collection tubes. Additives in the tubes preserve the sample, and the order of draw minimizes contamination between tubes.
\section*{History}
The practice of bloodletting by venesection dates to antiquity, but modern venepuncture for diagnostic testing developed with laboratory medicine in the 19th and 20th centuries, replacing crude lancets with the vacuum tube system introduced in the 1940s.
\section*{Why It Is Done}
Ordered to obtain blood for complete blood counts, metabolic panels, cultures, serology, coagulation studies, therapeutic drug monitoring, and countless other laboratory assays, and to screen, diagnose, or monitor disease.
\section*{Is This a Primary or Secondary Test or Procedure}
A primary procedure whose purpose is specimen collection; it is supportive rather than therapeutic and is the gateway to laboratory testing.
\section*{How It Is Performed}
The patient is seated or supine with the arm extended. A tourniquet is applied, and a suitable vein in the antecubital fossa is selected by inspection and palpation. The site is cleaned with alcohol, the vein is anchored, and the needle is advanced at a shallow angle. Blood is drawn into evacuated tubes in the correct order, the tourniquet is released, and the needle is withdrawn with pressure applied to the site.
\section*{Preparation}
Informed consent and, for some tests, fasting. The collector should confirm patient identification, test requirements, and collection tubes before the draw.
\section*{Contraindications and Precautions}
Avoid the limb with an arteriovenous fistula, lymphedema, infection, or a previous mastectomy on that side. Use caution with anticoagulated patients, coagulopathies, and children, and avoid repeated probing of a site.
\section*{Risks}
Pain, bruising, hematoma, bleeding, nerve injury from deep probing, arterial puncture, infection, and vasovagal reactions or fainting. Hemolyzed or clotted samples may require repeat collection.
\section*{Aftercare and Recovery}
Pressure is held over the site for a few minutes and a small dressing is applied. The samples are labeled, transported to the laboratory, and analyzed according to the ordered tests.
\section*{Outcome and Follow-up}
Results depend on the tests ordered and are interpreted by the clinician in the context of the patient's history and examination. The quality of collection directly affects result reliability.
\section*{Expected Results}
Reference ranges vary by test, age, and sex; normal results fall within the laboratory's published reference intervals for the ordered panel.
\section*{Follow-up}
Abnormal results may prompt repeat testing, additional confirmatory or reflex testing, and further investigations directed by the clinical picture.
\section*{When to Seek Medical Care}
Follow the procedure team's specific instructions. Seek urgent medical assessment for severe or worsening pain, uncontrolled bleeding, fever or signs of infection, trouble breathing, chest pain, fainting, new weakness or confusion, inability to urinate, or any rapidly worsening symptom. The appropriate warning signs depend on the procedure and the patient's medical condition.

\section*{References}
\begin{enumerate}
  \item MedlinePlus. Venipuncture. U.S. National Library of Medicine; 2025.
  \item Current Medical Diagnosis and Treatment. McGraw-Hill; 2025.
\end{enumerate}

\clearpage
